% =====================================================================
%  16  結言 — 成果と課題・展望の1枚サマリ
%  source_memo: 全章＋trial #201。13_conclusion.tex は本文未記載のため、
%   緒言〜分析の確定値から要約を構成。型K。
%  方針：上=1文サマリ、左=成果、右=課題と展望。色は赤字の数値のみ。丸括弧禁止。
% =====================================================================
\begin{frame}[t]

  \vspace{0.4em}
  {\bfseries\large LPG から PDH でプロピレン $402$ kt/年・$99.5$ wt\% を製造する全プロセスを設計し、全 $23$ 変数の同時最適化で TAC を最小化した}\par

  \vspace{0.7em}

  \begin{columns}[T,onlytextwidth]
    % ===== 左：成果 =============================================
    \begin{column}{0.5\textwidth}
      {\bfseries 成果}\par
      \vspace{0.3em}
      {\footnotesize\raggedright
        \textbf{反応器}　浅床軸流・多基スイングと床内 HGM 補温で断熱平衡の上限を緩和し、単通転化率 $40.6$ \%・選択率 $80.1$ \%\par
        \vspace{0.4em}
        \textbf{分離}　脱ブタン・脱エタン・C3 スプリッタの $3$ 塔に、最難の $\mathrm{C_3H_6}/\mathrm{C_3H_8}$ を膜で粗濃縮するハイブリッドを併用、副生水素は PSA で $99.9$ mol\% 超に回収\par
        \vspace{0.4em}
        \textbf{最適化}　全 $23$ 変数を TPE ベイズ最適化で同時決定、熱統合で用役費を $214.7$ 億円/年削減\par}
    \end{column}

    % ===== 右：課題と展望 =======================================
    \begin{column}{0.5\textwidth}
      {\bfseries 課題と展望}\par
      \vspace{0.3em}
      {\footnotesize\raggedright
        反応速度論に内在する選択率の頭打ちで総合収率は $81.2$ \% に留まり、難分離 C3 に設備費が集中する\par
        \vspace{0.4em}
        原料費が TAC の $65.9$ \% を占め、現行想定 $95$ 円/kg では\,{\boldmath\color{SpRed}赤字 $223.7$ 億円/年}\par
        \vspace{0.4em}
        損益分岐は LPG $64.4$ 円/kg\par
        $\to$\,\textbf{安価な原料を確保できる立地なら採算は大きく改善し得る}\par}
    \end{column}
  \end{columns}

  \vfill

  {\setlength{\fboxsep}{6pt}\setlength{\fboxrule}{1pt}%
   \fcolorbox{HeaderBlue}{white}{%
    \begin{minipage}{\dimexpr\linewidth-16pt}
     \footnotesize\centering
     反応・分離・水素回収・熱統合を貫く全体同時最適化により、物理に根ざした最良設計点 trial \#201 を提示した。
    \end{minipage}}}

\end{frame}
